% !TeX root = ../../Book3.tex
\section{亚纯延拓}
\label{b3:sec:2102}

递推式既能向右计算，也能向左延拓。后一用法必须检查不同移动次数给出的值相同，并辨明分母的零点是真极点还是可去点。

\subsection{递推关系}
\label{b3:subsec:210201}

对任意正整数 \(N\)，在 \(\operatorname{Re}z>-N\) 定义
\[
\Gamma(z)=\frac{\Gamma(z+N)}{z(z+1)\cdots(z+N-1)}.
\]
当分母不为零时这是全纯函数，其他点至多为极点。用一次递推消去 \(z+N\) 可知，采用 \(N+1\) 与 \(N\) 的定义在重叠区域一致。因此这些局部表达式粘合成整个复平面上的亚纯函数，并继续满足递推式。以后 \(\Gamma\) 均指这一延拓。

\subsection{极点}
\label{b3:subsec:210202}

\begin{proposition}\label{b3:prop:21-poles}
Gamma 函数的极点恰为 \(0,-1,-2,\ldots\)，每个都是一阶极点。
\end{proposition}
\begin{proof}
若 \(z\) 不是非正整数，选足够大的 \(N\) 可使分子处于右半平面，且分母各因子非零，故该点全纯。在 \(z=-m\) 附近取 \(N=m+1\)。分子在该点等于 \(\Gamma(1)=1\)，分母只有因子 \(z+m\) 为零，其余因子都非零，所以是恰好一阶的极点。
\end{proof}
延拓并不赋予原积分在 \(\operatorname{Re}z\le0\) 上收敛的含义。例如 \(\Gamma(-1/2)=-2\sqrt\pi\)，而形式积分 \(\int_0^\infty e^{-t}t^{-3/2}\,\mathrm{d}t\) 发散且被积函数为正；两者不能混为一谈。

\subsection{留数}
\label{b3:subsec:210203}

上述局部表示直接给出
\[
\operatorname{Res}_{z=-m}\Gamma(z)=\frac{(-1)^m}{m!},\quad m=0,1,2,\ldots.
\]
事实上，提出 \(z+m\) 后，其他因子在 \(-m\) 的乘积为 \((-1)^m m!\)。在反演积分中移动围道时，这串留数将产生幂级数的系数，而极点位置给出相应的幂次。

\ProseParagraphBreak

复共轭关系 \(\Gamma(\overline z)=\overline{\Gamma(z)}\) 先由右半平面的积分得到，再经恒等定理延拓。负实轴相邻极点之间的符号因此可用递推明确计算：在 \((-m-1,-m)\) 上符号为 \((-1)^{m+1}\)。此符号变化来自越过一阶极点，不意味着存在实零点；事实上下一节将证明复零点也不存在。

\ProseParagraphBreak

极点附近还能得到有限部分，而不仅是留数。先由乘积的首项展开得到 \(\Gamma(z)=z^{-1}-\gamma+O(z)\)。在 \(z=-m+\varepsilon\) 使用递推，分母非零部分为
\[
 \prod_{j=0}^{m-1}(-m+\varepsilon+j)
 =(-1)^m m!\left(1-H_m\varepsilon+O(\varepsilon^2)\right).
\]
分子 \(\Gamma(1+\varepsilon)=1-\gamma\varepsilon+O(\varepsilon^2)\)。因此
\[
 \Gamma(-m+\varepsilon)=\frac{(-1)^m}{m!}
       \left(\frac1\varepsilon+H_m-\gamma+O(\varepsilon)\right),
 \qquad H_0=0.
\]
当把 Gamma 与另一个含零点的函数相乘时，留数决定可去后的值，有限部分则决定导数。只知道极点阶数不足以完成这些局部计算。
\ProseParagraphBreak
参数正则化提供一个与递推不同的延拓图像。在 \(-1<\operatorname{Re}z<0\) 中，积分
\[
 \int_0^1(e^{-t}-1)t^{z-1}\,\mathrm{d}t
 +\int_1^\infty e^{-t}t^{z-1}\,\mathrm{d}t+\frac1z
\]
绝对收敛并等于 Gamma 延拓。因为 \(e^{-t}-1=O(t)\)，第一项比原 Euler 积分多出一个可积幂次；最后的 \(1/z\) 恰把减掉的常数项补回。继续减去更多 Taylor 项，就能向更左的半平面推进。所谓解析延拓在这里表现为一项明确的奇部加一个普通收敛积分，并没有把发散积分本身赋予任意的数值。
